\section*{Synthèse pour la direction}
\addcontentsline{toc}{section}{Synthèse pour la direction}

MEDICA-SERVICES héberge et exploite les applications de gestion de douze établissements de soins privés \piece{(P01)}. Son système d'information porte des identifiants de patients et, pour une partie des clients, des données de santé \piece{(P05)}. Une altération non détectée ou une indisponibilité prolongée peut donc atteindre la prise en charge des patients chez des tiers. Ces enjeux conduisent l'autodiagnostic DCSSI à une somme de \SommeMaxima\ points, soit un \textbf{niveau adéquat de maturité de \Cible}, sans ajustement proposé. Ce niveau exprime la rigueur souhaitable et non la sécurité déjà atteinte.

Les pièces démontrent une maturité effective de \textbf{niveau 1} pour la stratégie, la gestion des risques, la supervision et la réalisation des mesures, et de \textbf{niveau 2} pour les règles, la conception et l'exploitation. Aucun processus n'atteint le niveau 3, car aucune pratique n'est standardisée sur tout le périmètre. L'écart à la cible est donc de deux à trois niveaux. Deux déclarations de la direction ne sont pas étayées par les pièces techniques. La DSI affirme que tous les comptes privilégiés ont un MFA alors que l'export n'en montre que six sur dix enrôlés \piece{(T04, P17)}, ce qui la contredit. La DG affirme que toutes les restaurations sont testées chaque trimestre alors que deux services sur neuf seulement disposent d'un procès-verbal \piece{(T09, P17)}, ce qui ne prouve pas l'absence des autres tests mais laisse l'affirmation non démontrée.

Trois risques dominent l'analyse. Le premier est la prise de contrôle d'un compte à privilèges. Le 9 octobre 2026, le compte admin-old, sans MFA et dont la dernière activité documentée remontait à février 2025 \piece{(T03)}, s'est connecté au VPN après trois échecs puis a modifié un groupe d'administrateurs. L'alerte AL-1048 s'est ensuite fermée automatiquement, sans ticket ni examen tracé \piece{(T07, T08)}. Ces signaux appellent une investigation, sans constituer à eux seuls la preuve d'une intrusion. Le deuxième risque est l'impossibilité de reprendre le service dans les 8 heures contractuelles après un rançongiciel, puisque les postes peuvent joindre le réseau des sauvegardes sans restriction \piece{(FW-31, T02)} et que les restaurations sont peu éprouvées. Le troisième concerne l'application métier APP01, à laquelle manque un correctif critique \piece{(T05)} sur un hôte en retard de correctifs de 74 jours sans exception validée \piece{(T06)}. Son score est inférieur à ceux de R04 et R06, mais nous le plaçons parmi les risques majeurs pour son impact maximal sur les données de patients et parce qu'il se réduit rapidement.

Nous recommandons à la direction de prendre trois décisions. Elle peut d'abord mandater dès la remise de ce rapport l'investigation d'AL-1048, avec des mesures conservatoires sur admin-old et la restriction en régie des règles FW-31 et FW-17. Elle peut ensuite engager \TotalRetenu\ sur le plafond de \Plafond, réparti entre six actions qui couvrent les comptes à privilèges, la supervision, les correctifs, la restauration, la gestion des risques et la révision de la PSSI. Une marge de \Marge\ reste disponible. Il lui revient enfin d'accepter formellement, pour douze mois, les risques résiduels liés aux mesures différées, au premier rang desquelles la segmentation complète et la standardisation des projets. Cette trajectoire vise un niveau 2 démontré sur les sept processus à douze mois, première étape cumulative vers la cible.
